\section{从运动 BCI 到语音 BCI}

\subsection{运动 BCI 的早期成功}

基于对运动皮层神经编码的理解，研究者们成功构建了多种运动 BCI 系统。

\begin{figure}[H]
\centering
\includegraphics[width=0.95\textwidth]{slides-images/slide-016.jpg}
\caption{BCI 系统示意图：脑内微电极阵列记录神经信号，通过线缆传输到解码计算机，解码运动意图\protect\footnotemark}
\end{figure}
\footnotetext{Slides 第17页。红色问号标示脊髓损伤位置，黄色线路为 BCI 绕行通路。}

\subsubsection{2D 光标控制与虚拟键盘}

NPTL 实验室在 2017 年展示了一项重要成果：瘫痪患者可以通过脑机接口控制虚拟键盘上的光标进行打字，平均速度约为\textbf{每分钟 20 个正确字符}，峰值可达 40 个字符/分钟。这意味着患者不再需要依赖他人帮助字母板拼读，可以\textbf{独立}通过思维控制计算机进行沟通。

\subsubsection{机械臂控制}

另一个标志性成果是利用 BCI 控制机械臂。Caltech 的研究展示了瘫痪患者通过大脑信号控制机械臂拿起饮料并送到嘴边饮用——这是患者多年来第一次能够自己完成这样的动作。

\subsubsection{手写 BCI}

2021 年，NPTL 实验室的 Frank Willett 发表论文，展示了\textbf{手写 BCI}：通过解码患者``想象书写''时的运动皮层信号，可以将手写字母转化为文本，速度达到约\textbf{每分钟 18 个词}，远快于 2D 光标方式的 8 个词/分钟。

\begin{importantbox}{不同沟通方式的速度对比}
\begin{itemize}
    \item \textbf{Sip-and-puff 接口}：$\sim$5 词/分钟
    \item \textbf{2D 光标 BCI}：$\sim$8 词/分钟
    \item \textbf{手写 BCI}：$\sim$18 词/分钟
    \item \textbf{自然说话}：$\sim$150--160 词/分钟
\end{itemize}
即使最好的运动 BCI 也远未达到自然沟通的速度。问题是：\textbf{能否直接恢复语音}？
\end{importantbox}

\subsection{为什么语音 BCI 比运动 BCI 更难}

语言的产生是一个极其复杂的神经过程，涉及大脑多个区域的协调：

\begin{itemize}
    \item \textbf{知识与推理区域}（右侧）：负责语义理解和高层认知
    \item \textbf{语义和句法区域}（中央）：负责词语选择和语法组织
    \item \textbf{语音感知区域}（左侧）：负责言语感知
    \item \textbf{运动规划与执行区域}：负责控制口腔面部肌肉产生语音
\end{itemize}

语音产生涉及的肌肉运动比手臂运动\textbf{复杂得多、快速得多}——嘴唇、舌头、声带、下颌等多个发声器官（articulators）需要在毫秒级时间内精确协调。

\begin{warningbox}{语音 BCI 的特殊挑战}
\begin{itemize}
    \item 语音涉及多个发声器官的\textbf{快速协调运动}，远比手臂运动复杂
    \item 对于丧失语言能力的患者，无法直接测量其发声器官的运动（因为发声器官本身也可能瘫痪）
    \item 大脑中与语言相关的区域广泛分布，不像手臂运动那样集中在运动皮层的特定区域
    \item 训练数据采集困难——患者无法正常说话，数据量必然有限
\end{itemize}
\end{warningbox}

\subsection{关键突破：从发声器官运动到音素解码}

为了降低问题的复杂度，研究者们没有试图直接解码每个发声器官的连续运动轨迹，而是转向解码\textbf{离散的音素（phoneme）}。

\begin{knowledgebox}{音素：语言的基本单位}
\textbf{音素（phoneme）}是一种语言中能区分词义的最小语音单位。英语中共有约 \textbf{39--44 个音素}，包括元音（如 /a/、/i/）和辅音（如 /p/、/t/、/s/）。每个音素对应特定的发声器官配置——舌头位置、嘴唇形状、声带振动方式等。

选择音素作为解码目标的优势：
\begin{itemize}
    \item 音素集合小（$\sim$40 个），远少于词汇量（数万个），数据需求少
    \item 不同音素在运动皮层中具有\textbf{可区分的}神经活动模式
    \item 从音素可以组合出任意词汇，不受词汇表限制
\end{itemize}
\end{knowledgebox}

2021 年，UCSF 的研究者们使用\textbf{皮层电图（ECoG）}技术（电极放在皮层表面、不穿入大脑）首次证明了语音 BCI 的可行性：能够以约 75\% 的准确率解码 50 个词的小词汇表。虽然词汇量有限，但这标志着语音 BCI 从概念走向了现实。

\subsection{本章小结}

\begin{itemize}
    \item 运动 BCI 已经成功帮助瘫痪患者控制光标、机械臂和进行手写输入
    \item 但运动 BCI 的沟通速度远低于自然说话速度
    \item 语音 BCI 是下一个目标，但面临更大的技术挑战
    \item 解码离散音素（而非连续运动轨迹）是降低问题复杂度的关键策略
    \item UCSF 2021 年首次证明了小词汇量语音 BCI 的可行性
\end{itemize}

%% ========== Part 4: 高性能语音神经假体 ==========

